%% kapa: exam=1 round=28 year=2017 period=1 form=A subject=economics subject_ko=경제학원론 seq=21 number=61
%% source: Q-Net 감정평가사 기출문제 (cst00309) artlSeq=5199603 file=extracted/1008/5199603/2017년 제28회 감정평가사 1차시험 문제지/1교시 A형.pdf
%% conversion: kordoc 4.24.51 -> tools/gen1.py (auto); review=proofread-sonnet; flags=none
\begin{question}{61}
실업에 관한 설명으로 옳지 않은 것은?
\begin{choices}
  \item 일자리를 가지고 있지 않으나 취업할 의사가 없는 사람은 경제활동인구에 포함되지 않는다.
  \item 실업이란 사람들이 일할 능력과 의사를 가지고 일자리를 찾고 있으나 일자리를 얻지 못한 상태를 말한다.
  \item 자연실업률은 구조적 실업만이 존재하는 실업률이다.
  \item 실업자가 구직을 단념하여 비경제활동인구로 전환되면 실업률이 감소한다.
  \item 경기변동 때문에 발생하는 실업은 경기적(cyclical) 실업이다.
\end{choices}
\end{question}
\answer{3}
